\PassOptionsToPackage{numbers,sort&compress}{natbib}
\documentclass{article}

\usepackage[final]{neurips_2025}

\makeatletter
\renewcommand{\@notice}{}
\makeatother

\usepackage[utf8]{inputenc}
\usepackage[T1]{fontenc}
\usepackage[english]{babel}

\usepackage{microtype}
\usepackage{amsmath,amssymb,amsfonts,bm}
\usepackage{booktabs}
\usepackage{multirow}
\usepackage{array}
\usepackage{graphicx}
\usepackage{subcaption}
\usepackage{url}
\usepackage{hyperref}
\usepackage{natbib}
\usepackage{float}
\usepackage{placeins}

\renewcommand{\topfraction}{0.9}
\renewcommand{\bottomfraction}{0.8}
\renewcommand{\textfraction}{0.07}
\renewcommand{\floatpagefraction}{0.8}

\captionsetup{justification=centering}

% % Reduce espacio alrededor de figuras y tablas (sin romper NeurIPS)
% \setlength{\textfloatsep}{10pt plus 2pt minus 2pt}
% \setlength{\floatsep}{8pt plus 2pt minus 2pt}
% \setlength{\intextsep}{10pt plus 2pt minus 2pt}

% % Reduce espacio de captions
% \setlength{\abovecaptionskip}{5pt}
% \setlength{\belowcaptionskip}{0pt}

\title{Does Graph Geometry Help Forecasting the S\&P 100?}

\author{
Antonio Lorenzo Díaz Meco and Javier Ríos Montes \\
Final Project --- Geometric Deep Learning \\
\href{https://github.com/Javirios03/sp100-graph-forecasting}{GitHub Repository}
}

\begin{document}

\maketitle

\begin{abstract}
    Write the abstract here.
\end{abstract}

\section{Introduction}

Stock forecasting is a challenging task that, due to its economic implications, has been an attractive area of research for decades. In specific, AI and machine learning techniques have been widely applied to this problem, with varying degrees of success. Despite the advances in the matter, many current approaches are limited to using technical analysis to capture historical trends of each stock price. In this paper, we explore the potential benefits of incorporating graph geometry into stock forecasting models, specifically focusing on the S\&P 100 index. We decide to focus on this subset of the stock market to reduce the complexity of the problem and to be able to focus on the relationships between the stocks, which is where we expect graph geometry to provide the most value.

In indeces such as the S\&P 100, there are important relationships between individual stocks, such as shared market sectors or financial indicators, that may not be fully captured by traditional forecasting models and that could be better modeled using graph-based approaches. By representing the stocks as nodes in a graph and their relationships as edges, we can leverage geometric deep learning techniques to capture these interactions and, potentially, improve forecasting performance

In order to evaluate the effectiveness of graph geometry in stock forecasting, we first implement two baseline models: a Random Forest and an LSTM. These models will serve as benchmarks to compare against a Graph Neural Network (GNN). The choice of these architectures is motivated by their widespread use in time series forecasting and their ability to capture different aspects of the data. The Random Forest is a powerful ensemble method that can capture non-linear relationships (and, given a sufficiently informative feature set, obtain performance comparable to more complex models), while the LSTM is a type of recurrent neural network that is particularly well-suited for sequential data (in this case, the temporal evolution of the index). The GNN, on the other hand, is designed to capture the relationships between stocks by modeling them as a graph, which is where we expect to see the most significant improvements in forecasting performance.

We model the relationships between stocks using a graph structure, where nodes represent individual stocks and edges represent relationships, mainly shared market sectors, correlations between stock returns, and divergences in the returns' distributions. The problem is thus framed as a node classification task, in which we aim to predict the direction of the stock price movement (up, neutral, or down) for each stock in the index within a given time frame. We evaluate the performance of the models using standard metrics such as accuracy, precision, recall, and F1-score, and we analyze the results to understand the impact of incorporating graph geometry into the forecasting process.

\section{Related Work}

\section{Background}

\section{Method}

\section{Experimental Setup}

\subsection{Data and Preprocessing}

\subsection{Models and Hyperparameters}

\section{Results and Discussion}

\section{Conclusions and Future Work}

\newpage

\section*{Contribution Statement}

While both of us developed the code together and contributed to the writing of the report, there are certain sections that were implemented primarily by each one of us.

\begin{itemize}
    \item \textbf{Antonio Lorenzo Díaz Meco}: Code for the 3 models (Random Forest, LSTM, and GNN), and experimental setup. Writing of the Method and Experimental Setup.
    \item \textbf{Javier Ríos Montes}: Code for data preprocessing and Repository structure. Writing of the Introduction, Related Work, Background, Results and Discussion, and Conclusions and Future Work.
\end{itemize}

\bibliographystyle{plainnat}
\bibliography{references}

\end{document}
